% chapters/01_introduction.tex

\section{Overview}

This chapter establishes the intellectual foundation for the thesis. It begins with a
concrete financial-market narrative that motivates mean-field games with common noise
(\Cref{sec:what-is-mfg}), derives the governing forward-backward stochastic differential
equation (FBSDE) system from first principles (\Cref{sec:fbsde}), and surveys the
literature across stochastic partial differential equation (SPDE) theory, reinforcement
learning, and quantitative finance (\Cref{sec:state-of-art}). We identify three precise
gaps that this work addresses, explicitly situating each within the existing body of
knowledge. The six principal contributions are then presented together with a candid
discussion of their novelty relative to prior work and their interconnections. We
demonstrate the deep structural connections between the three fields, provide a
ten-chapter roadmap, and collect all key notation and standing assumptions. Three
tailored reading paths are offered, and the chapter closes with a discussion of the
mathematical, computational, and market conditions that make this research both
possible and timely.

By the end of this chapter the reader will be able to answer five questions:

\begin{enumerate}
    \item How does common noise fundamentally alter the geometry and stability of
    mean-field equilibria?
    \item What is the FBSDE system, and what role does each equation play in
    characterising the equilibrium?
    \item What does the literature already contain, and what are the three specific
    gaps this thesis addresses --- understood as incremental extensions rather than
    fundamental breakthroughs?
    \item What exactly does this thesis prove, and how do the six principal results
    interconnect and relate to prior work?
    \item How should one navigate the remaining nine chapters?
\end{enumerate}

\paragraph{Executive summary.} This thesis models financial markets as a large
population of interacting institutional investors. Each investor adjusts its portfolio
(defined as log-positions across assets) based on the aggregate behaviour of all
others --- the distribution of positions across the market --- which evolves as a
stochastic process taking values in the space of probability measures equipped with the
Wasserstein-2 geometry. The Wasserstein-2 geometry is the central analytical and
economic object of the thesis, providing both the metric for stability analysis and the
basis for measuring market inefficiency. An equilibrium is reached when the distribution
generated by everyone acting optimally coincides with the distribution that each
investor takes as given.

The central claim of this thesis is that \textbf{market inefficiency can be productively
understood as distributional disequilibrium under shared shocks} --- a measurable
distance in probability space. This thesis reframes market inefficiency as a geometric
misalignment in the space of distributions, measured by the Wasserstein-2 distance
between the empirical market measure and the theoretical equilibrium measure.

A key structural result is that common noise preserves the contraction structure of the
MFG fixed-point map in expectation under synchronous coupling and homogeneous
common-noise exposure. The contraction is established at the level of expected
Wasserstein dynamics; it does not extend to pathwise contraction or heterogeneous noise
structures. The first-order stochastic terms cancel in expectation, and the contribution
of the common noise enters only through quadratic variation terms in the It\^o expansion
of optimally coupled processes. This cancellation is a second-moment phenomenon
formalised in Theorem~B.9.

We prove that this model is mathematically stable (in the sense of a contraction in a
weighted Wasserstein metric) even under common shocks, design regularised reinforcement
learning algorithms with convergence guarantees (established for regularised, compact
function classes as a theoretical baseline; extending these results to overparameterised
deep neural networks remains an open problem), and show that deviations from equilibrium
yield economically meaningful trading signals whose directional behaviour aligns with
the theoretical structure.

This thesis develops a mathematically rigorous and computationally implementable
framework for mean-field games with common noise. It integrates three research streams
that have not yet been combined into a single framework simultaneously providing
(i) quantitative stability estimates, (ii) learning guarantees under common noise, and
(iii) empirical validation on real financial data:

\begin{itemize}
    \item \textbf{SPDE theory} --- the analysis of forward-backward stochastic partial
    differential equations arising from the Stochastic Maximum Principle;
    \item \textbf{Reinforcement learning} --- multi-timescale actor-critic algorithms
    with provable convergence guarantees for regularised function classes;
    \item \textbf{Quantitative finance} --- deployable alpha generation from
    equilibrium theory, validated by rigorous statistical testing.
\end{itemize}

The framework yields six principal contributions spanning well-posedness, numerical
analysis, learning theory, and empirical finance. Selected core estimates are supported
by formal verification in Lean 4, and the empirical application provides evidence broadly
consistent with the equilibrium predictions. \textbf{All theoretical results are subject
to the standing assumptions collected in \Cref{sec:notation-assumptions}; these
assumptions are strong and their implications for practical implementation are discussed
throughout.}


\section{What Is a Mean-Field Game?}
\label{sec:what-is-mfg}

\subsection{The $N$-Player Problem: A Financial Market}

Consider a financial market populated by $N$ institutional investors --- hedge funds,
asset managers, pension funds, and proprietary trading desks. We work on a complete
filtered probability space $(\Omega, \mathcal{F}, (\mathcal{F}_t)_{t \in [0,T]},
\mathbb{P})$ satisfying the usual conditions. The filtration $(\mathcal{F}_t)$ is the
augmentation of that generated by a collection of independent Brownian motions: a
$d_0$-dimensional common noise $W^0 = (W_t^0)_{t \in [0,T]}$ and, for each agent
$i \in \{1, \dots, N\}$, a $d$-dimensional idiosyncratic noise $W^i = (W_t^i)_{t \in [0,T]}$,
with $W^i \perp\!\!\!\perp W^j$ for $i \neq j$ and all $W^i$ independent of $W^0$. All
stochastic processes are assumed to be progressively measurable with respect to this
filtration. The set of admissible controls $\mathcal{A}_{\mathrm{adm}}$ consists of
$\mathcal{F}_t$-adapted processes with values in a compact convex set
$A \subset \mathbb{R}^k$ satisfying $\mathbb{E}\int_0^T |\alpha_t|^2\, dt < \infty$.

Each investor $i$ controls its own portfolio state $X_t^i \in \mathbb{R}^d$, representing
the vector of log-positions across $d$ assets (i.e., the natural logarithm of the
fraction of wealth allocated to each asset, or equivalently the log-portfolio weights).
The objective is to maximise a risk-adjusted return over the finite horizon $[0,T]$:
\begin{equation}
    J_i(\alpha^i, \alpha^{-i}) = \mathbb{E}\left[ \int_0^T f\big(t, X_t^i, \mu_t^N,
    \alpha_t^i\big)\, dt + g\big(X_T^i, \mu_T^N\big) \right],
    \label{eq:objective}
\end{equation}
where
\begin{itemize}
    \item $\alpha_t^i \in A \subset \mathbb{R}^k$ is investor $i$'s control (trading
    rate), assumed to be $\mathcal{F}_t$-adapted;
    \item $\alpha^{-i}$ denotes the strategies of all other investors;
    \item $f : [0,T] \times \mathbb{R}^d \times \mathcal{P}_2(\mathbb{R}^d) \times A \to
    \mathbb{R}$ is the running reward (return net of trading costs);
    \item $g : \mathbb{R}^d \times \mathcal{P}_2(\mathbb{R}^d) \to \mathbb{R}$ is the
    terminal reward (e.g., liquidation value).
\end{itemize}

The payoff depends on the empirical measure of the population,
\begin{equation}
    \mu_t^N := \frac{1}{N} \sum_{j=1}^{N} \delta_{X_t^j} \in \mathcal{P}_2(\mathbb{R}^d),
    \label{eq:empirical-measure}
\end{equation}
which captures price impact, crowding effects, and information spillovers. Here
$\mathcal{P}_2(\mathbb{R}^d)$ denotes the space of probability measures on
$\mathbb{R}^d$ with finite second moment, equipped with the Wasserstein-2 distance
$W_2$. The dependence of $J_i$ on $\alpha^{-i}$ enters solely through the empirical
measure $\mu_t^N$, reflecting the mean-field interaction.

Each portfolio evolves according to the McKean--Vlasov stochastic differential equation
(SDE)
\begin{equation}
    dX_t^i = b\big(t, X_t^i, \mu_t^N, \alpha_t^i\big)\, dt + \sigma\, dW_t^i +
    \sigma_0\, dW_t^0, \qquad X_0^i \sim \mu_{\mathrm{init}},
    \label{eq:state-sde}
\end{equation}
where
\begin{itemize}
    \item $b : [0,T] \times \mathbb{R}^d \times \mathcal{P}_2(\mathbb{R}^d) \times A \to
    \mathbb{R}^d$ is the drift (expected price movement net of trading impact);
    \item $\sigma > 0$ is the idiosyncratic volatility; for $d > 1$ assets we set
    $\sigma$ to be a scalar (for notational simplicity);
    \item $\sigma_0 \in \mathbb{R}^{d \times d_0}$ is the common-noise coefficient,
    and $W_t^0 \in \mathbb{R}^{d_0}$ is a shared Brownian motion representing
    aggregate shocks;
    \item $\mu_{\mathrm{init}} \in \mathcal{P}_2(\mathbb{R}^d)$ is the initial
    distribution of positions.
\end{itemize}

\begin{example}[Linear-quadratic portfolio MFG]
\label{ex:lq-mfg}
For a single asset ($d=1$), set
\[
    b(t,x,\mu,\alpha) = \kappa(\bar{x} - x) + \alpha, \qquad
    f(t,x,\mu,\alpha) = -\left(x^2 + \frac{\lambda}{2}\alpha^2\right), \qquad
    g(x,\mu) = -(x - \bar{x}_T)^2,
\]
where $\bar{x} = \int y\, \mu(dy)$ is the population mean. Here the investor seeks to
minimise a quadratic cost functional. This canonical LQ-MFG admits an explicit solution
via Riccati equations and serves as the computational benchmark throughout Chapters 4
and 9.
\end{example}

\subsection{Why the $N$-Player Problem Is Intractable}

The $N$-player game is a stochastic differential game whose Nash equilibrium requires
solving $N$ coupled Hamilton--Jacobi--Bellman equations on $\mathbb{R}^{dN}$.
Fundamental obstacles include dimensionality (for $d=50$, $N=200$ the state space has
dimension 10{,}000), broken symmetry (heterogeneity in realised paths destroys
closed-form solutions except in LQ cases), and exponential complexity of dynamic
programming. These difficulties motivate the mean-field limit.

\subsection{The Mean-Field Limit: $N \to \infty$}

\paragraph{The Key Observation.} Lasry and Lions (2007) and independently Huang,
Malham\'e, and Caines (2006) observed that as $N \to \infty$ the empirical measure
$\mu_t^N$ concentrates on a (possibly random) measure flow $\mu_t$. By the conditional
law of large numbers given the common noise $W^0$, and using conditional propagation of
chaos,
\[
    \mu_t^N \xrightarrow[N \to \infty]{W_2} \mu_t,
\]
with convergence rate $O(N^{-1/2})$ for sufficiently regular measures (Fournier and
Guillin, 2015; see Chapter 2). Each investor then solves a single-agent control problem
against the aggregate measure flow $\mu = (\mu_t)_{t \in [0,T]}$.

\begin{definition}[MFG Nash equilibrium]
A pair $(\mu^*, \alpha^*)$ is an MFG Nash equilibrium if:
\begin{enumerate}[label=(\arabic*)]
    \item \textit{Optimality}: Given $\mu^*$, the control $\alpha^*$ maximises the
    representative agent's objective with the empirical measure replaced by $\mu^*$;
    \item \textit{Consistency}: For almost all $t \in [0,T]$,
    $\mu_t^* = \mathcal{L}(X_t^* \mid \mathcal{F}_t^{W^0})$, where
    $\mathcal{L}(X_t^* \mid \mathcal{F}_t^{W^0})$ is a random probability measure
    measurable with respect to $\mathcal{F}_t^{W^0}$, and $X^*$ satisfies
    \eqref{eq:state-sde} under $\alpha^*$ with $\mu^N$ replaced by $\mu^*$.
\end{enumerate}
\end{definition}

This is the fixed-point problem $\Phi(\mu) = \mu$ on the space of measure flows, where
$\Phi(\mu) := \mathcal{L}(X^\mu \mid \mathcal{F}_T^{W^0})$. \textbf{All contraction
statements in this thesis refer to the fixed-point map on measure flows, not to
pathwise contraction of individual trajectories.} Theorem~B.9 establishes that, under
monotonicity and sufficient regularity ensuring stability of optimal transport maps,
$\Phi$ is a contraction locally in time in general, and globally under strengthened
convexity conditions in the weighted Wasserstein-2 metric
$\|\mu^1 - \mu^2\|_{\lambda, T} := \sup_{t \leq T} e^{-\lambda t}
\mathbb{E}[W_2(\mu_t^1, \mu_t^2)^2]^{1/2}$, guaranteeing existence and uniqueness.

\paragraph{The $\varepsilon$-Nash Property.} If all $N$ investors adopt $\alpha^*$,
the resulting profile is an $\varepsilon_N$-Nash equilibrium with
$\varepsilon_N = O(N^{-1/2})$ (Carmona and Delarue, 2018). Chapter 4 provides a
quantitative bound calibrating this error against realistic transaction costs.

\subsection{Why Common Noise Is Non-Trivial}

\paragraph{Deterministic Case ($\sigma_0 = 0$).} The measure flow $\mu_t$ is
deterministic and the MFG reduces to a coupled system of PDEs. Existence, uniqueness,
and stability follow under monotonicity conditions (Lasry and Lions, 2007; Carmona and
Delarue, 2018).

\paragraph{Common-Noise Case ($\sigma_0 > 0$): New Challenges.} When $\sigma_0 > 0$,
the measure flow $\mu_t$ becomes a random element of $\mathcal{P}_2(\mathbb{R}^d)$,
evolving as a stochastic process on an infinite-dimensional space. Three fundamental
difficulties appear: (i) the Fokker--Planck equation becomes an SPDE; (ii) the
consistency condition is measure-valued and must hold conditional on
$\mathcal{F}^{W^0}$; (iii) proving convergence of particle systems requires compactness
criteria on the Skorokhod space $D([0,T]; \mathcal{P}_2(\mathbb{R}^d))$ of c\`adl\`ag
measure-valued processes, endowed with the usual $J_1$ topology (Jakubowski's criterion
or Mitoma's theorem; see Chapter 2).

\paragraph{Key Structural Observation.} The key asymmetry is that individual states are
adapted to the full filtration $\mathcal{F}_t$, while the measure flow is adapted only
to the common-noise filtration $\mathcal{F}_t^{W^0}$. When two candidate MFG solutions
$(\mu^1, \mu^2)$ are coupled on the same probability space sharing the identical common
noise $W^0$, the common-noise terms partially cancel in expectation under synchronous
coupling and homogeneous common-noise exposure. \textbf{The contraction is established
at the level of expected Wasserstein dynamics; it does not imply pathwise
cancellation.} The first-order stochastic differential cancels only after taking
expectations, and the contribution of the common noise enters through Lipschitz
estimates on the adjoint BSDE and the control difference, affecting the drift of the
expected squared Wasserstein distance. This cancellation is a second-moment
(in-expectation) phenomenon under synchronous coupling, formalised in Theorem~B.9.


\section{The FBSDE Characterisation}
\label{sec:fbsde}

\subsection{Overview of the System}

The MFG equilibrium of \Cref{sec:what-is-mfg} can be characterised by a coupled
forward-backward stochastic differential equation (FBSDE) system, coupled with a
measure-valued SPDE --- the central mathematical object of this thesis. Every
subsequent chapter either analyses this system theoretically (Chapters 2--3),
discretises it (Chapter 4), learns it via reinforcement learning (Chapters 5--6), or
extracts tradable signals from it (Chapters 7--9).

\subsection{Heuristic Derivation}

The dynamic programming principle yields a backward SPDE for the value function; the
stochastic maximum principle gives the adjoint BSDE; the forward population evolution
is governed by the Fokker--Planck SPDE; consistency closes the loop. Full rigorous
derivation is provided in Chapter 2, following the probabilistic approach of Carmona
and Delarue (2018).

\subsection{The Forward Equation: Population Measure Flow}

\begin{definition}[Fokker--Planck SPDE]
Under the MFG equilibrium, the population measure $\mu_t \in \mathcal{P}_2(\mathbb{R}^d)$
satisfies, for all test functions $\phi \in C_c^\infty(\mathbb{R}^d)$,
\begin{equation}
    d\langle \mu_t, \phi \rangle = \langle \mu_t, \mathcal{L}_t^* \phi \rangle\, dt +
    \langle \mu_t, \nabla \phi \cdot \sigma_0 \rangle\, dW_t^0, \qquad
    \mu_0 = \mu_{\mathrm{init}},
    \label{eq:fokker-planck}
\end{equation}
where $\langle \mu, \phi \rangle := \int \phi\, d\mu$ and $\mathcal{L}_t^*$ is the
formal adjoint of the generator.
\end{definition}

The rigorous interpretation uses a Gelfand triple $V \subset H \subset V^*$ with
$H = L^2(\mathbb{R}^d)$, $V = H^1(\mathbb{R}^d)$; solutions live in the Skorokhod space
$D([0,T]; \mathcal{P}_2(\mathbb{R}^d))$ (see Chapter 2).

\subsection{The Backward Equation: Adjoint Variables}

\begin{definition}[Adjoint BSDE]
The adjoint process $(p_t, q_t, r_t)$ satisfies
\begin{equation}
    -dp_t = \partial_x H(t, X_t, \mu_t, p_t, \alpha_t^*)\, dt - q_t\, dW_t^i -
    r_t\, dW_t^0, \qquad p_T = \partial_x g(X_T, \mu_T),
    \label{eq:adjoint-bsde}
\end{equation}
with Hamiltonian $H = p \cdot b + f$.
\end{definition}

All processes are adapted to $\mathcal{F}_t$, while $\mu_t$ is adapted to
$\mathcal{F}_t^{W^0}$. Under sufficient smoothness (Assumption A3), formal
identification with Lions derivatives holds; these identifications are not required for
existence proofs.

\subsection{The Optimality Condition}

\begin{proposition}[Stochastic Maximum Principle]
Under Assumption A5 (uniform concavity of $\alpha \mapsto H$), the optimal control
satisfies $\alpha_t^* = \arg\max_{\alpha \in A} H(t, X_t, \mu_t, p_t, \alpha)$. For the
LQ example, $\alpha_t^* = p_t / \lambda$.
\end{proposition}

\subsection{The Consistency Condition}

\begin{equation}
    \mu_t = \mathcal{L}\big(X_t^* \mid \mathcal{F}_t^{W^0}\big) \qquad
    \text{for all } t \in [0,T],
    \label{eq:consistency}
\end{equation}
where $X^*$ follows \eqref{eq:state-sde} with $\mu^N$ replaced by $\mu$ and control
$\alpha^*$. This is proved simultaneously with existence in Theorem~B.9.

\subsection{The Complete FBSDE System}

The system is solved as a fixed point in the space of measure flows: for a given $\mu$,
solve the forward--backward system for $(X, p)$, update $\mu$ via consistency, and
iterate. The five coupled components are:
\begin{align}
    \text{(Fwd)} \quad & dX_t^* = b(t, X_t^*, \mu_t, \alpha_t^*)\, dt + \sigma\, dW_t^i
        + \sigma_0\, dW_t^0, \quad X_0^* \sim \mu_{\mathrm{init}}, \notag \\
    \text{(FP)}  \quad & d\langle \mu_t, \phi \rangle = \langle \mu_t,
        \mathcal{L}_t^* \phi \rangle\, dt + \langle \mu_t, \nabla \phi \cdot
        \sigma_0 \rangle\, dW_t^0, \quad \forall \phi \in C_c^\infty(\mathbb{R}^d), \notag \\
    \text{(Bwd)} \quad & -dp_t = \partial_x H(t, X_t, \mu_t, p_t, \alpha_t^*)\, dt -
        q_t\, dW_t^i - r_t\, dW_t^0, \quad p_T = \partial_x g(X_T, \mu_T),
        \label{eq:fbsde-system} \\
    \text{(Opt)} \quad & \alpha_t^* = \argmax_{\alpha \in A} H(t, X_t, \mu_t, p_t,
        \alpha), \notag \\
    \text{(Cons)} \quad & \mu_t = \mathcal{L}(X_t^* \mid \mathcal{F}_t^{W^0}). \notag
\end{align}

Theorem~B.9 establishes, under Assumptions A1--A7, existence, uniqueness, and a
quantitative stability estimate. Each component corresponds to an observable or
estimable market quantity, enabling mapping to tradable signals (Chapter 8).


\section{State of the Art}
\label{sec:state-of-art}

This section surveys the literature across three domains: mean-field game theory,
reinforcement learning for MFGs, and quantitative finance applications. We highlight
both foundational results and recent advances, and we identify the specific gaps that
this thesis addresses. Throughout, we emphasise that this thesis builds incrementally on
recent work, particularly Carmona and Delarue (2018), Ahuja, Ren, and Yang (2019),
Angiuli et al. (2024), and Cardaliaguet and Lehalle (2018).

\subsection{Mean-Field Game Theory}

\paragraph{Foundations.} The mean-field game framework was introduced independently by
Lasry and Lions (2007) and Huang, Malham\'e, and Caines (2006). Lasry and Lions
established existence and uniqueness of solutions for deterministic MFGs under
monotonicity conditions via PDE methods. The probabilistic formulation was developed by
Carmona and Delarue (2018), who provided a comprehensive treatment using forward-backward
SDEs and extended the theory to the common-noise setting. Cardaliaguet et al. (2019)
introduced the master equation, a single equation on the Wasserstein space that encodes
the entire MFG equilibrium, and developed the Lions derivative calculus necessary for
its analysis.

\paragraph{Common-noise MFGs.} The presence of common noise introduces substantial
mathematical challenges. Ahuja, Ren, and Yang (2019) proved well-posedness for a class
of common-noise MFGs with linear-convex structure and weak-monotone cost functions,
using a continuation method for FBSDEs with monotone functionals. Cardaliaguet, Seeger,
and Souganidis (2025) established SPDE regularity results for MFGs with common noise
and degenerate idiosyncratic noise. \textbf{The contraction method used in this thesis
builds directly on the probabilistic FBSDE framework of Carmona and Delarue (2018) and
the monotone functional approach of Ahuja et al. (2019). The incremental contribution
is the derivation of explicit, parameter-dependent stability constants that are not
readily available in the prior literature.}

\paragraph{Trade crowding.} Cardaliaguet and Lehalle (2018) formulated the classical
optimal liquidation problem as a mean-field game of controls, providing a closed-form
solution and demonstrating that traders do not need instantaneous knowledge of the
whole system but can learn by observing others' behaviour. Lehalle and Mouzouni (2019)
extended this framework to portfolios of correlated instruments, deriving optimal
trading flows and analysing their impact on perceived correlations using real data on
176 US stocks. Neuman and Vo{\ss} (2023) studied a multi-player liquidation game with
common signals and proved convergence at rate $O(N^{-2})$. This thesis builds upon this
line of work by incorporating common noise and deriving a richer set of trading signals,
most notably the novel crowding measure $c_t = W_2(\mu_t, \mu_t^*)$.

\subsection{Reinforcement Learning for Mean-Field Games}

\paragraph{Two-timescale methods.} The application of reinforcement learning to MFGs
has seen rapid progress. Angiuli et al. (2022) introduced a unified two-timescale RL
algorithm that solves either MFG or mean-field control problems depending on the ratio
of learning rates. Hu and Lauri\`ere (2023) developed a two-timescale actor-critic
method specifically for common-noise MFGs. These works demonstrate convergence for
finite state-action spaces under discretisation, but do not provide explicit
Wasserstein control or extend to the continuous-time SPDE setting.

\paragraph{Three-timescale methods and common noise.} Most relevant to this thesis,
Angiuli, Fouque, Lauri\`ere, and Zhang (2024) proved the convergence of a
three-timescale Q-learning algorithm for mean-field control games with common noise,
establishing that the multi-scale dynamics correctly capture the representative agent's
viewpoint. Their proof generalises Borkar's two-timescale approach to three timescales,
but the setting is restricted to finite state and action spaces. \textbf{The present
thesis extends this ODE-method approach to the continuous-time, continuous-space
setting with explicit Wasserstein distance control. The proof structure follows that of
Angiuli et al. (2024) closely; the novel technical challenge is establishing tightness
and convergence of the measure-valued component in the Wasserstein topology.}

\paragraph{Entropy regularisation and exponential convergence.} Recent work has
established exponential convergence rates for entropy-regularised RL. Ma, Li, and
Zhang (2024) proved that entropy-regularised policy gradient achieves exponential
convergence under a log-Sobolev inequality. Theorem 5.2 of this thesis adapts this
approach to the MFG setting, providing an exponential rate for the Wasserstein distance
between learned and equilibrium measures. \textbf{The result holds under an idealised
log-Sobolev assumption and serves as a theoretical benchmark rather than a guarantee for
practical deep RL implementations.}

\subsection{Quantitative Finance}

\paragraph{Factor models and crowding.} Classical factor models (Fama and French, 1993)
lack equilibrium grounding. The crowding problem --- the tendency of quantitative
strategies to become overcrowded and subsequently underperform --- has been documented
by Khandani and Lo (2011), but lacked a theoretically derived measure of crowding. The
optimal execution literature (Almgren and Chriss, 2000) is predominantly single-agent.

\paragraph{MFG applications in finance.} Beyond the trade crowding models of Cardaliaguet
and Lehalle (2018) and Lehalle and Mouzouni (2019), several works have applied MFGs to
financial problems. Carmona, Fouque, and Sun (2015) modelled systemic risk as a
mean-field game of interbank borrowing and lending. Casgrain and Jaimungal (2020)
studied algorithmic trading with price impact using MFG methods. This thesis
contributes to this growing literature by providing a unified framework from MFG
equilibrium to empirically validated trading signals, with a novel emphasis on
distributional disequilibrium measured by the Wasserstein distance.

\subsection{Wasserstein Distance in Financial Applications}

The use of the Wasserstein metric in finance has gained traction. Fournier and Guillin
(2015) established the rate of convergence of the empirical measure in Wasserstein
distance, providing the theoretical foundation for particle approximations. Esfahani
and Kuhn (2018) developed distributionally robust optimisation using Wasserstein balls,
providing tractable reformulations and performance guarantees. Blanchet, Kang, and
Murthy (2019) applied Wasserstein profile inference to machine learning problems. Kim,
Korajczyk, and Neuhierl (2015) showed that temporal differences in return distributions
measured by Wasserstein distance predict cross-sectional returns. The present thesis
uses the Wasserstein distance in a novel way: as a measure of deviation from the MFG
equilibrium, capturing distributional disequilibrium anchored to an economic
equilibrium concept.

\subsection{Gaps in the Literature}

The survey reveals three specific gaps, each framed as an incremental extension of
existing work:

\begin{itemize}
    \item \textbf{Gap 1 (Theoretical)}: While existence and uniqueness of common-noise
    MFGs have been established (Ahuja et al., 2019; Carmona and Delarue, 2018), explicit
    parameter-dependent stability bounds suitable for numerical analysis and learning
    algorithms are not readily available. This thesis contributes Theorem~B.9, which
    provides such constants under additional regularity assumptions.

    \item \textbf{Gap 2 (Algorithmic)}: Existing RL frameworks for MFGs either assume
    finite state-action spaces (Angiuli et al., 2024) or lack explicit Wasserstein
    control. No approach simultaneously provides convergence in the common-noise,
    continuous-space setting, explicit Wasserstein control, and a multi-timescale
    architecture aligned with the SPDE structure. This thesis contributes Theorems
    5.1--5.2 and the hierarchical architecture of Chapter 6, albeit under strong
    function-approximation assumptions.

    \item \textbf{Gap 3 (Applied)}: A fully integrated pipeline combining MFG
    equilibrium solution, theoretically justified trading signals, and rigorous
    empirical validation on real equity data remains largely unexplored. The work of
    Lehalle and Mouzouni (2019) provides an empirical analysis of correlations but does
    not derive and test a full trading strategy. This thesis resolves this gap through
    the alpha signals of Chapter 8 and the empirical backtest of Chapter 9.
\end{itemize}


\section{The Three Critical Gaps}

\paragraph{Gap 1 (Theoretical).} Explicit parameter-dependent stability bounds suitable
for numerical and learning algorithms are not readily available. While Ahuja, Ren, and
Yang (2019) and Carmona and Delarue (2018) established well-posedness, they did not
provide the quantitative contraction constants needed for algorithmic design. Addressed
by Theorem~B.9.

\paragraph{Gap 2 (Algorithmic).} No existing RL framework simultaneously provides
provable convergence to the equilibrium measure flow in the common-noise,
continuous-space setting, with explicit Wasserstein control and a multi-timescale
architecture aligned with the SPDE structure. Angiuli et al. (2024) proved
three-timescale convergence for finite spaces; this thesis extends the ideas to
continuous spaces under compactness and Lipschitz assumptions. \textbf{The extension is
not trivial, but the proof methodology follows the same ODE framework.} Addressed by
Theorems 5.1--5.2 and Proposition 6.1.

\paragraph{Gap 3 (Applied).} A fully integrated pipeline combining MFG equilibrium
solution, theoretically justified trading signals, and empirical validation on real
equity data remains largely unexplored. Resolved by the sequence: Chapter 3 (theory)
$\to$ Chapter 4 (numerics) $\to$ Chapter 8 (signals) $\to$ Chapter 9 (validation).


\section{Principal Contributions}

All results are subject to the standing assumptions of \Cref{sec:notation-assumptions}
and additional structural conditions discussed in corresponding chapters. Each
contribution is explicitly situated relative to prior work. \textbf{Theorem~B.9 is the
central structural result; all subsequent results --- numerical stability, RL
convergence, and empirical signal construction --- depend explicitly on its contraction
estimate.}

\subsection*{Theory (Well-Posedness)}

\begin{theorem}[Quantitative well-posedness]
Under Assumptions A1--A7, the FBSDE system \eqref{eq:fbsde-system} admits a unique
solution. Under monotonicity and sufficient regularity ensuring stability of optimal
transport maps (e.g., absolute continuity with uniformly bounded densities), the
solution map is a contraction (locally in time in general, globally under strengthened
convexity) in the weighted Wasserstein-2 metric, yielding an explicit stability
estimate
\[
    \mathbb{E}\big[W_2(\mu_t^1, \mu_t^2)^2\big] \leq C_0 \int_0^t e^{C(t-s)}
    \mathbb{E}\big[W_2(\mu_s^1, \mu_s^2)^2\big]\, ds,
\]
where $C = 2L_b^2 + 2\|\sigma_0\|_F^2 + C_{\mathrm{reg}}$, with $C_{\mathrm{reg}}$
depending on bounds of Lions derivatives of $b, f, g$. The $\varepsilon$-Nash property
holds with $\varepsilon_N = O(N^{-1/2})$. (Chapter 3)
\end{theorem}

\noindent\textit{Relation to prior work:} Extends Carmona--Delarue (2018) and Ahuja et
al. (2019) by providing explicit parameter dependence in the contraction constant. The
contraction is in expectation, not pathwise.

\subsection*{Numerics}

\begin{theorem}[Numerical convergence]
For the Euler--Maruyama discretisation of the Fokker--Planck SPDE, under sufficient
smoothness, the weak error is $O(\Delta t^{1/2} + \Delta x)$. A CFL condition is
derived. (Chapter 4)
\end{theorem}

\noindent\textit{Relation to prior work:} Builds on Chassagneux--Crisan--Delarue (2019)
and Achdou--Porretta (2016), extending to the common-noise SPDE case.

\subsection*{Learning}

\begin{theorem}[Convergence of regularised multi-timescale actor-critic]
A regularised multi-timescale actor-critic algorithm with learning rates satisfying
Robbins--Monro and timescale separation, under bounded Lipschitz function approximators
(e.g., regularised neural networks), converges $\mathbb{P}$-almost surely along
subsequences (in the sense of $\liminf W_2 = 0$) to a fixed point of the limiting
mean-field operator. Under uniqueness (Theorem~B.9), the limit is the MFG equilibrium.
(Chapter 5)
\end{theorem}

\noindent\textit{Relation to prior work:} Extends Angiuli et al. (2024) from finite
spaces to continuous spaces with Wasserstein control; the subsequential convergence is
a weaker result than full almost-sure convergence, which remains open.

\begin{theorem}[Exponential rate under entropy regularisation]
Under entropy regularisation and a log-Sobolev inequality for the policy class (a
restrictive assumption illustrating achievable rates under idealised conditions), the
policy iteration scheme satisfies
\[
    W_2(\mu_t, \mu^*) \leq C_0 e^{-\lambda t} W_2(\mu_0, \mu^*), \qquad
    \lambda = \frac{\beta}{2 C_{\mathrm{Lip}}}.
\]
(Chapter 5)
\end{theorem}

\noindent\textit{Relation to prior work:} Adapts Ma, Li, and Zhang (2024) to the MFG
setting; the LSI assumption limits practical applicability.

\begin{proposition}[Hierarchical architecture]
The SPDE decomposition implies a natural timescale separation exploited by a
hierarchical RL architecture (SAC/PPO/TD3). (Chapter 6)
\end{proposition}

\noindent\textit{Relation to prior work:} Provides a structural motivation for
algorithm choice; the specific algorithm mapping is heuristic, not proven optimal.

\subsection*{Finance}

\begin{proposition}[Four alpha signals]
The equilibrium structure implies four candidate signals: optimal spread,
mean-reversion speed, factor exposure, and crowding measure
$c_t = W_2(\mu_t, \mu_t^*)$. The optimal spread, mean-reversion speed, and factor
exposure are adaptations of signals found in Cardaliaguet--Lehalle (2018) and
Lehalle--Mouzouni (2019); the crowding measure $c_t$ is the novel contribution. Unlike
prices, $\mu_t^*$ is not directly observable; estimation introduces model risk
(addressed in Chapter 9). Mean-variance portfolio weights with crowding discount are
provided. (Chapter 8)
\end{proposition}

\subsection*{Empirical Finding}

\begin{theorem}[Empirical validation]
A strategy based on Proposition~8.1 signals, applied to S\&P 500 constituents
(2000--2024, survivorship-bias-free), yields an estimated annualised Sharpe ratio of
1.9 under the model and estimation procedure, after transaction costs (5--10 bps) and
slippage. Statistical significance: $p = 0.002$ (10{,}000 block-bootstrap). Maximum
drawdown $-18\%$. Metrics use purged walk-forward cross-validation. Robustness checks
support qualitative conclusions. \textbf{The reported performance is conditional on
model specification and estimation of $\mu_t^*$; the magnitude of the Sharpe ratio is
sensitive to these choices. Model risk dominates statistical risk}, and the empirical
results should be interpreted as validation of the structural predictions of the model
rather than evidence of a directly deployable trading strategy. The strategy assumes a
small-player limit; a capacity analysis is provided with strong caveats. (Chapter 9)
\end{theorem}

\paragraph{Formal verification.} The Gronwall estimate of Theorem~B.9, the Riccati
system derivation, the CFL condition of Theorem 1.7, and the optimality of the linear
feedback control have been formally verified in Lean 4. The code is available in the
supplementary repository and is excerpted in Appendix~B. \textbf{The verification is
partial; full verification of the entire pipeline remains future work.}

\begin{table}[ht]
\centering
\caption{Summary of Principal Contributions and Their Novelty}
\label{tab:contributions}
\begin{tabular}{p{3cm}p{4cm}p{5cm}}
\toprule
\textbf{Contribution} & \textbf{Prior Work} & \textbf{Incremental Novelty} \\
\midrule
Theorem 3.1 (Quantitative well-posedness) & Carmona--Delarue (2018), Ahuja et al. (2019)
    & Explicit parameter-dependent contraction constant; contraction in expectation \\
\addlinespace
Theorem 4.1 (Numerical convergence) & Chassagneux--Crisan--Delarue (2019),
    Achdou--Porretta (2016) & Combined scheme with CFL condition for common-noise SPDE \\
\addlinespace
Theorem 5.1 (RL convergence) & Angiuli et al. (2024) & Continuous-space extension with
    Wasserstein control; subsequential convergence \\
\addlinespace
Theorem 5.2 (Exponential rate) & Ma--Li--Zhang (2024) & Application to MFG; idealised
    LSI benchmark \\
\addlinespace
Proposition 6.1 (Hierarchical architecture) & Hu--Lauri\`ere (2023) & SPDE-grounded
    timescale decomposition; heuristic SAC-PPO-TD3 mapping \\
\addlinespace
Proposition 8.1 (Alpha signals) & Cardaliaguet--Lehalle (2018), Lehalle--Mouzouni (2019)
    & Four-signal framework; novel crowding measure $c_t$ \\
\addlinespace
Empirical Finding 9.1 & Lehalle--Mouzouni (2019) & Full out-of-sample strategy on
    S\&P 500; Sharpe 1.91 (conditional on model) \\
\bottomrule
\end{tabular}
\end{table}


\section{Connections Between Fields}

\paragraph{Connection 1 (SPDE $\to$ RL architecture).} The master equation's
decomposition into idiosyncratic, mean-field, and common-noise operators motivates a
hierarchical RL architecture (SAC/PPO/TD3). Proposition~E.2 provides timescale
separation.

\paragraph{Connection 2 (Stability $\to$ Crisis robustness).} The coupled Wasserstein
dynamics show the contraction constant remains controlled under large common shocks in
expectation. The crowding signal $c_t = W_2(\mu_t, \mu_t^*)$ provides a quantitative
measure of market stress. (In extreme stress, orthogonality properties may degrade;
this is acknowledged as a limitation.)

\paragraph{Connection 3 (Equilibrium $\to$ Alpha).} Definition~7.1 formalises market
efficiency as $\mu_t = \mu_t^*$. The Wasserstein distance measures deviation from Nash
equilibrium. When $c_t$ is large, profitable deviations are predicted; when small, the
market is approximately efficient.


\section{Thesis Roadmap}

\subsection{Chapter Dependency Graph}

Theory stream: Chapter 2 $\to$ 3 $\to$ 4 $\to$ 5 $\to$ 6. Finance stream: Chapter 7
$\to$ 8 $\to$ 9. All synthesised in Chapter 10.

\subsection{Chapter-by-Chapter Summary}

\begin{itemize}
    \item \textbf{Chapter 2}: Preliminaries (Wasserstein geometry, It\^o calculus,
    SPDEs, Lions derivative).
    \item \textbf{Chapter 3}: Quantitative well-posedness (Theorem~B.9).
    \item \textbf{Chapter 4}: Numerical analysis (Theorem 1.7).
    \item \textbf{Chapter 5}: RL convergence (Theorems D.1, D.5).
    \item \textbf{Chapter 6}: Hierarchical architecture (Proposition E.2).
    \item \textbf{Chapter 7}: Case study --- systemic risk, bridging theory to
    application.
    \item \textbf{Chapter 8}: Alpha generation (Proposition 1.11).
    \item \textbf{Chapter 9}: Empirical validation (Empirical Finding 1.12).
    \item \textbf{Chapter 10}: Synthesis, limitations, open problems.
\end{itemize}


\section{Key Notation and Standing Assumptions}
\label{sec:notation-assumptions}

\textit{Reference section; may be skipped on first reading.}

\subsection{Notation Reference}

\begin{center}
\begin{tabular}{ll}
\toprule
\textbf{Symbol} & \textbf{Meaning} \\
\midrule
$\mathcal{P}_2(\mathbb{R}^d)$ & probability measures with finite second moment \\
$W_2(\mu, \nu)$ & Wasserstein-2 distance \\
$\langle \mu, \phi \rangle$ & $\int \phi\, d\mu$ \\
$\mathcal{L}(X)$ & law of $X$ \\
$W_t^i, W_t^0$ & idiosyncratic and common Brownian motions \\
$\mathcal{F}_t^{W^0}$ & common-noise filtration \\
$b, f, g$ & drift, running reward, terminal reward \\
$H = p \cdot b + f$ & Hamiltonian \\
$p_t, q_t, r_t$ & adjoint processes \\
$U(t, \mu)$ & master-equation value function; $D_\mu U$: Lions derivative \\
$L_b, L_f, L_g$ & Lipschitz constants \\
$\beta$ & entropy regularisation coefficient \\
$c_t = W_2(\mu_t, \mu_t^*)$ & crowding measure \\
$D([0,T]; E)$ & Skorokhod space (c\`adl\`ag), $J_1$ topology \\
$C_c^\infty(\mathbb{R}^d)$ & smooth compactly supported functions \\
$\|\sigma_0\|_F$ & Frobenius norm \\
\bottomrule
\end{tabular}
\end{center}

\subsection{Standing Assumptions}

\begin{assumption}[Lipschitz drift]
$\|b(t,x,\mu,\alpha) - b(t,x',\mu',\alpha')\| \leq L_b(\|x-x'\| + W_2(\mu,\mu') +
\|\alpha - \alpha'\|)$.
\end{assumption}

\begin{assumption}[Lipschitz rewards]
$f$ and $g$ are Lipschitz.
\end{assumption}

\begin{assumption}[Additional regularity]
$b, f, g$ are $C^1$ in $x$ and Lions-differentiable in $\mu$, with bounded first
derivatives (and second derivatives when required).
\end{assumption}

\begin{assumption}[Bounded common noise]
$\|\sigma_0\|_F \leq M_\sigma < \infty$.
\end{assumption}

\begin{assumption}[Uniform concavity of Hamiltonian]
$H(t,x,\mu,p,\alpha') \leq H(t,x,\mu,p,\alpha) + \partial_\alpha H \cdot (\alpha' -
\alpha) - \frac{\lambda_H}{2}\|\alpha' - \alpha\|^2$.
\end{assumption}

\begin{assumption}[Regular terminal reward]
$g$ is $C^1$ in $x$ and Lions-differentiable in $\mu$, with bounded derivatives.
\end{assumption}

\begin{assumption}[Initial condition]
$\mu_{\mathrm{init}}$ has finite fourth moment.
\end{assumption}

\begin{remark}
Assumptions A1--A7 are standard in the MFG literature but are restrictive. In
particular, A3 (Lions differentiability) and the requirement of absolute continuity for
optimal transport stability are strong and are not satisfied by empirical discrete
measures without regularisation. The theoretical results should be viewed as
establishing a well-posedness baseline under idealised conditions; practical
implementation requires regularisation and model risk quantification.
\end{remark}


\section{Guide to Reading}

\begin{itemize}
    \item \textbf{Path A (Pure mathematics)}: Chapters 2, 3, 5.
    \item \textbf{Path B (Algorithms/RL)}: Chapters 2 (overview), 5, 6.
    \item \textbf{Path C (Quantitative finance)}: Chapters 2 (Wasserstein), 7, 8, 9.
\end{itemize}


\section{Why This Problem, Why Now}

Three developments converged in 2023--2025: mathematical maturity (Cardaliaguet, Seeger,
and Souganidis 2025 resolved SPDE regularity for common-noise MFGs; Ahuja, Ren, and Yang
2019 established well-posedness for monotone functionals), computational feasibility
(GPU-accelerated optimal transport via the Sinkhorn algorithm of Cuturi (2013), mature
RL libraries), and market structure (post-COVID concentrated participation, common
shocks, crowding events).

\paragraph{Central message.} This thesis argues that modern financial markets are
naturally modelled as mean-field games: large populations of agents optimise against
aggregate behaviour under common shocks. A key structural result is that common noise
preserves the contraction structure of the MFG fixed-point map in expectation under
synchronous coupling and homogeneous exposure. The Wasserstein distance
$W_2(\mu_t, \mu_t^*)$ provides a mathematically coherent measure of crowding ---
sensitive to displacement and mass redistribution, not merely density differences. When
this distance is large, theory predicts profitable deviations; when small, the market
is approximately efficient. Proposition~1.11 turns this insight into tradable signals;
Chapter 9 provides empirical evidence, with appropriate statistical caution, that the
approach has out-of-sample merit. \textbf{In this sense, the thesis reframes market
inefficiency not as a price deviation, but as a distributional disequilibrium under
shared shocks.}
